% ===============================================================
%  جلسه بیست‌وچهارم
% ===============================================================
\chapter{بازه اطمینان با نامساوی چبی‌شف و مقدمه تخمین بیزی}

این جلسه دو بخش دارد. در بخش اول، تخمین بازه‌ای جلسه قبل را از توزیع گاوسی به \textbf{توزیع‌های دلخواه} گسترش می‌دهیم: وقتی شکل توزیع را نمی‌دانیم و فقط واریانس (یا کرانی برای آن) معلوم است، نامساوی چبی‌شف یک بازه اطمینان \textbf{تضمین‌شده} می‌دهد. در بخش دوم، دیدگاه کاملاً تازه‌ای را آغاز می‌کنیم: \textbf{تخمین بیزی}، که در آن علاوه بر داده‌ها، \textbf{اطلاعات پیشین} درباره پارامتر هم در اختیار داریم و برای به کار بردن آن، خود پارامتر را یک متغیر تصادفی مدل می‌کنیم. این بخش با معرفی تخمین‌گر \textbf{بیشینه توزیع پسین \lr{(MAP)}} تمام می‌شود.

\section{تخمین بازه‌ای میانگین برای توزیع‌های دلخواه}

\subsection{صورت مسئله}
فرض کنید متغیر تصادفی $X$ میانگین نامعلوم $\mu$ دارد و درباره شکل توزیعش چیزی نمی‌دانیم (گاوسی، یکنواخت، دوکوهانه، \ldots). تنها اطلاع ما درباره واریانس به یکی از دو شکل زیر است:
\begin{itemize}
\item مقدار دقیق واریانس معلوم و برابر $\sigma^2$ است، یا
\item مقدار دقیق واریانس معلوم نیست، اما می‌دانیم عدد معلوم $\sigma^2$ یک \textbf{کران بالا} برای واریانس است: $\Var X\le\sigma^2$.
\end{itemize}
حالت دوم در عمل بسیار رایج است. مثلاً اگر $X$ همواره در بازه $[a,b]$ باشد، بدون دانستن توزیع می‌دانیم $\Var X\le\frac{(b-a)^2}{4}$ (برای متغیر برنولی: $\Var X=p(1-p)\le\frac14$). اگر $n$ نمونه \lr{i.i.d.} از $X$ داشته باشیم، می‌خواهیم یک تخمین‌گر بازه‌ای $(\theta_l,\theta_u)$ برای میانگین، با سطح اطمینان دلخواه $1-\delta$ بسازیم.

\subsection{استفاده از نامساوی چبی‌شف}
یک راه مستقیم و سرراست برای یافتن یک تخمین‌گر بازه‌ای \textbf{نسبتاً مناسب}، استفاده از نامساوی چبی‌شف است. سه واقعیت را کنار هم می‌گذاریم:

\textbf{۱) $\bar X_n$ در هر حالتی بی‌گرایش است.} چون امید ریاضی خطی است: $\E\bar X_n=\frac1n\sum_i\E X_i=\mu$؛ این به شکل توزیع ربطی ندارد.

\textbf{۲) واریانس $\bar X_n$ (یا کران بالای آن) برابر $\sigma^2/n$ است.} چرا؟ چون نمونه‌ها مستقل‌اند، واریانس مجموع برابر مجموع واریانس‌هاست:
\[
\Var\bar X_n=\Var\Big(\frac1n\sum_{i=1}^nX_i\Big)=\frac{1}{n^2}\sum_{i=1}^n\Var X_i\ \le\ \frac{1}{n^2}\cdot n\sigma^2=\frac{\sigma^2}{n}
\]
(اگر واریانس دقیقاً $\sigma^2$ باشد، نامساوی به تساوی تبدیل می‌شود).

\textbf{۳) نامساوی چبی‌شف} (جلسه ۱۸): برای هر متغیر تصادفی $Y$ با واریانس متناهی و هر $\varepsilon>0$:
\[
\Prob\big(|Y-\E Y|\ge\varepsilon\big)\le\frac{\Var Y}{\varepsilon^2}.
\]
این نامساوی به \textbf{هیچ} اطلاعاتی درباره شکل توزیع نیاز ندارد.

اکنون چبی‌شف را برای $Y=\bar X_n$ و $\varepsilon=k\sigma/\sqrt n$ به کار می‌بریم:
\[
\Prob\Big(\big|\bar X_n-\mu\big|>\frac{k\sigma}{\sqrt n}\Big)\le\frac{\Var\bar X_n}{k^2\sigma^2/n}\le\frac{\sigma^2/n}{k^2\sigma^2/n}=\frac{1}{k^2}.
\]

\subsection{تعیین $k$ و ساخت بازه}
اگر سطح اطمینان \lr{(Level of Confidence)} را برابر $1-\delta$ بگیریم، کافی است $\frac1{k^2}=\delta$، یعنی $k=\frac{1}{\sqrt\delta}$ قرار دهیم:
\[
\Prob\Big(\big|\bar X_n-\mu\big|>\frac{k\sigma}{\sqrt n}\Big)\le\frac1{k^2}=\delta
\quad\Rightarrow\quad
\Prob\Big(\big|\bar X_n-\mu\big|\le\frac{\sigma}{\sqrt{n\delta}}\Big)\ge1-\delta .
\]
(رویداد دوم متمم رویداد اول است و $\frac{k\sigma}{\sqrt n}=\frac{\sigma}{\sqrt\delta\sqrt n}=\frac{\sigma}{\sqrt{n\delta}}$. در اسلاید علامت «$\ge=$» آمده که منظور همان «$\ge$» است.) مانند جلسه قبل، $|\bar X_n-\mu|\le d$ معادل $\bar X_n-d\le\mu\le\bar X_n+d$ است. پس:

\begin{theorem}[بازه اطمینان چبی‌شف برای میانگین]
اگر $X_1,\ldots,X_n$ نمونه‌های \lr{i.i.d.} از توزیعی دلخواه با میانگین $\mu$ و واریانس حداکثر $\sigma^2$ باشند، آنگاه تخمین‌گر بازه‌ای
\[
\Big(\ \bar X_n-\frac{\sigma}{\sqrt{n\delta}}\ ,\ \bar X_n+\frac{\sigma}{\sqrt{n\delta}}\ \Big)
\]
با احتمال \textbf{حداقل} $1-\delta$ میانگین واقعی $\mu$ را در بر می‌گیرد (یک جواب سازگار برای مسئله).
\end{theorem}

تفاوت مهم با جلسه قبل: آنجا احتمال پوشش \textbf{دقیقاً} $1-\delta$ بود (چون توزیع را دقیق می‌دانستیم)؛ اینجا فقط یک \textbf{کران پایین} داریم. احتمال واقعی پوشش ممکن است بسیار بیشتر از $1-\delta$ باشد، یعنی بازه بیش از حد لازم پهن است.

\subsection{ویژگی‌های این بازه}
طول بازه برابر $\dfrac{2\sigma}{\sqrt{n\delta}}$ است.
\begin{itemize}
\item \textbf{سازگاری:} با افزایش $n\to\infty$ طول بازه به سمت صفر میل می‌کند؛ این نشان از برقراری ویژگی \textbf{سازگاری \lr{(consistency)}} دارد.
\item \textbf{هزینه سنگین اطمینان:} با افزایش سطح اطمینان $1-\delta$ به سمت ۱، طول بازه به صورت $\delta^{-1/2}$ افزایش پیدا می‌کند که بسیار شدید است. دلیل این امر استفاده از نامساوی کلی و ضعیف چبی‌شف است که از بسیاری از ویژگی‌های تخمین‌گر $\bar X_n$ (مثلاً تقریباً گاوسی بودن آن طبق قضیه حد مرکزی) استفاده نمی‌کند.
\end{itemize}

جدول زیر ضریب $k$ (نسبت نصف طول بازه به $\sigma/\sqrt n$) را در دو روش مقایسه می‌کند:
\begin{center}
\renewcommand{\arraystretch}{1.4}
\begin{tabular}{c|c|c|c|c}
سطح اطمینان $1-\delta$ & $0.90$ & $0.95$ & $0.99$ & $0.999$\\\hline
چبی‌شف: $1/\sqrt\delta$ & $3.162$ & $4.472$ & $10.000$ & $31.623$\\
گاوسی: $z_{\delta/2}$ & $1.645$ & $1.960$ & $2.576$ & $3.291$\\\hline
نسبت طول دو بازه & $1.9$ & $2.3$ & $3.9$ & $9.6$
\end{tabular}
\end{center}

\begin{center}
\begin{tikzpicture}
\end{tikzpicture}\\
{\small ضریب $k$ برحسب سطح اطمینان: در روش چبی‌شف $k=1/\sqrt\delta$ با نزدیک شدن $1-\delta$ به ۱ بسیار سریع رشد می‌کند، اما ضریب گاوسی $z_{\delta/2}$ رشد بسیار کندی دارد (دم گاوسی نمایی کاهش می‌یابد).}
\end{center}

\begin{note}
\notetitle \textbf{چرا با این وجود از چبی‌شف استفاده می‌کنیم؟} چون \textbf{تضمین} می‌دهد. بازه گاوسی (یا بازه مبتنی بر قضیه حد مرکزی) فقط وقتی معتبر است که داده‌ها گاوسی باشند یا $n$ آن‌قدر بزرگ باشد که تقریب حد مرکزی خوب باشد؛ و اینکه «$n$ چقدر بزرگ کافی است» به شکل توزیع بستگی دارد که نمی‌دانیم. بازه چبی‌شف برای \textbf{هر} توزیع و \textbf{هر} $n$ (حتی $n=1$) درست است. این یک مصالحه کلاسیک است: \textbf{فرض کمتر $\Leftarrow$ نتیجه محافظه‌کارانه‌تر}.
\end{note}

\begin{note}
\notetitle \textbf{برای مطالعه بیشتر: نامساوی هوفدینگ.} اگر بدانیم $X\in[a,b]$ (کراندار)، نامساوی قوی‌تری وجود دارد:
\[
\Prob\big(|\bar X_n-\mu|\ge\varepsilon\big)\le2\exp\Big(-\frac{2n\varepsilon^2}{(b-a)^2}\Big),
\]
که با برابر $\delta$ قرار دادن سمت راست، نصف طول $\varepsilon=(b-a)\sqrt{\frac{\ln(2/\delta)}{2n}}$ به دست می‌آید. وابستگی به $\delta$ از $\delta^{-1/2}$ به $\sqrt{\ln(2/\delta)}$ کاهش یافته است، چون هوفدینگ از کرانداری (و در نتیجه همه گشتاورها) استفاده می‌کند، نه فقط از واریانس.
\end{note}

\begin{example}[بازه چبی‌شف برای نمره‌ها]
نمره دانشجویان یک درس عددی بین $0$ و $20$ است و توزیع آن نامعلوم است. نمره $n=100$ دانشجو به تصادف (با جایگذاری) انتخاب شده و میانگین $\bar x=13.2$ به دست آمده است. بازه اطمینانی با سطح حداقل ۹۵٪ برای میانگین نمره کل دانشجویان بیابید و آن را با بازه تقریبی مبتنی بر حد مرکزی مقایسه کنید.
\begin{solution}
\textbf{کران واریانس:} نمره در $[0,20]$ است، پس $\Var X\le\frac{(20-0)^2}{4}=100$، یعنی $\sigma=10$ یک کران بالای معتبر برای انحراف معیار است.

\textbf{بازه چبی‌شف:} $\delta=0.05$، پس
\[
\frac{\sigma}{\sqrt{n\delta}}=\frac{10}{\sqrt{100\times0.05}}=\frac{10}{\sqrt5}=4.47
\quad\Rightarrow\quad[13.2-4.47,\ 13.2+4.47]=[8.73,\ 17.67].
\]
\textbf{بازه تقریبی حد مرکزی} (با همان کران $\sigma=10$): $1.96\times\frac{10}{\sqrt{100}}=1.96$، پس $[11.24,\ 15.16]$.

بازه چبی‌شف حدود $2.3$ برابر پهن‌تر است، اما تضمین آن بدون هیچ فرض و تقریبی برقرار است.
\end{solution}
\end{example}

\section{تخمین بیزی \lr{(Bayesian Estimation)}}

\subsection{مفهوم تخمین بیزی}
\begin{itemize}
\item فرض کنید از توزیع پارامتریک $f_X(x;\theta)$ (دارای فضای پارامترهای $\Theta$) نمونه‌های $X_1,\ldots,X_n$ به صورت \lr{i.i.d.} حاصل شده‌اند.
\item هدف (در عمده موارد) یافتن یک تخمین نقطه‌ای از پارامتر نامعلوم $\theta$ است.
\item اما این بار بر خلاف موارد پیشین، مقداری \textbf{اطلاعات \lr{(information)}} در مورد پارامتر $\theta$ در دست است. یعنی تنها راه دسترسی ما به پارامتر(های) $\theta$ نمونه‌های $X_1,\ldots,X_n$ نیست، بلکه منبع اطلاعات دیگری نیز در کنار مشاهدات وجود دارد: \textbf{اطلاعات پیشین \lr{(Prior Information)}}.
\end{itemize}

در روش‌های قبلی (\lr{MLE}، روش گشتاورها، بازه اطمینان) فرض بود که پیش از دیدن داده‌ها \textbf{هیچ} چیزی درباره $\theta$ نمی‌دانیم و همه چیز را باید از داده‌ها استخراج کنیم. اما در دنیای واقعی تقریباً همیشه دانش قبلی داریم: احتمال شیر آمدن یک سکه معمولی نزدیک $\frac12$ است، قد انسان بالغ بین ۱ تا ۲٫۵ متر است، و \ldots

\begin{example}[انتخاب جراح]
فرض کنید قصد دارید برای انجام یک عمل جراحی خیلی مهم و حیاتی، بهترین پزشک را از میان تعدادی گزینه انتخاب کنید:
\begin{itemize}
\item برای هر یک از پزشکان، پرونده کاری چندساله آنان موجود است: اینکه هر کدام چه تعداد عمل جراحی انجام داده و چه تعدادی از آنان موفق بوده‌اند (\textbf{مشاهدات}).
\item همچنین، سابقه تحصیلات پزشکان نیز در دسترس است؛ یعنی می‌دانیم هر کدام از پزشکان فارغ‌التحصیل کدام دانشکده پزشکی بوده‌اند (\textbf{اطلاعات پیشین}).
\end{itemize}
این دو منبع اطلاعات را چگونه ترکیب کنیم؟
\begin{solution}
\textbf{تحلیل کیفی.} در صورتی که تعداد داده‌های موجود در سابقه کاری پزشکان کم باشد، به اطلاعات پیشین که مربوط به تحصیلات است اهمیت بیشتری داده خواهد شد. اما با بالا رفتن تعداد مشاهدات، اطلاعات پیشین به مرور رنگ باخته و درست‌نمایی بر پایه مشاهدات فعلی نقش پررنگ‌تری خواهد گرفت. به بیان دیگر، مصالحه (\lr{tradeoff}) بین اطلاعات پیشین و درست‌نمایی از روی دادگان با افزایش تعداد نمونه‌های موجود در داده تغییر خواهد کرد:
\begin{itemize}
\item در صورتی که شما تنها ۵ نمونه کار از پزشکان جراح را دیده باشید، احتمالاً تأکید بیشتری بر اطلاعات پیشین خواهید نمود.
\item اما در صورتی که تعداد نمونه‌کارهای ثبت‌شده افزایش یابد، دادگان به تدریج غنای کافی برای تخمین پارامتر را پیدا کرده و اثر اطلاعات پیشین به مرور کاهش می‌یابد.
\end{itemize}
\textbf{چرا \lr{MLE} اینجا ناکافی است؟} فرض کنید جراح «الف» ۴ عمل انجام داده که هر ۴ موفق بوده‌اند و جراح «ب» ۱۰۰ عمل با ۹۲ موفقیت. تخمین \lr{MLE} نرخ موفقیت (نسبت موفقیت‌ها) برای «الف» برابر $1$ و برای «ب» برابر $0.92$ است؛ پس \lr{MLE} جراح «الف» را انتخاب می‌کند! اما شهود ما می‌گوید ۴ موفقیت از ۴ عمل می‌تواند صرفاً خوش‌شانسی باشد و سابقه ۱۰۰ عمل «ب» قابل اعتمادتر است.

\textbf{اما چگونه می‌توان این فرآیند را کمّی نمود؟!} پاسخ، تخمین بیزی است که در ادامه می‌سازیم؛ حل کمّی همین مثال را در جلسه ۲۷ (توزیع‌های مزدوج) می‌بینیم.
\end{solution}
\end{example}

\subsection{مدل‌سازی تصادفی اطلاعات پیشین}
به روایتی، دو منبع اطلاعاتی که در مورد تخمین $\theta$ به ما کمک خواهند کرد، عبارتند از:
\begin{itemize}
\item \textbf{اطلاعات حاصل از نمونه‌های دیده‌شده} $X_1,\ldots,X_n$؛ به عبارت دیگر، \textbf{درست‌نمایی \lr{(Likelihood)}} مشاهدات به صورت تابعی از مقدار پارامتر.
\item \textbf{اطلاعات پیشین \lr{(Prior)}} داده‌شده در مورد پارامتر، که ارتباطی با مشاهدات فعلی ندارد.
\end{itemize}
درست‌نمایی مشاهدات به صورت تابعی از پارامتر را پیش‌تر توسط احتمال مشاهدات به شرط پارامتر مدل‌سازی کرده‌ایم: $\Prob_{X|\theta}\big(X_{1:n}\given\theta\big)$.

\textbf{ایده کلیدی:} اطلاعات پیشین در مورد پارامتر(ها) را نیز می‌توانیم به طریقی مشابه، از طریق یک توزیع مفروض و داده‌شده برای آنان مدل‌سازی نماییم: $\Prob_\theta(\theta)$. پس در این شیوه از تخمین، \textbf{خود پارامتر(ها) متغیر تصادفی هستند.}

\begin{note}
\notetitle \textbf{تصادفی بودن پارامتر یعنی چه؟} در دیدگاه بیزی، احتمال فقط «فراوانی در تکرار» نیست، بلکه \textbf{میزان باور} ما را هم توصیف می‌کند. وقتی می‌گوییم $\theta\sim f_\theta$، منظور این نیست که طبیعت واقعاً هر بار $\theta$ را از نو قرعه‌کشی می‌کند؛ منظور این است که \textbf{عدم قطعیت ما} درباره $\theta$ با این توزیع توصیف می‌شود. جایی که $f_\theta$ بزرگ است، مقادیری است که به نظر ما محتمل‌ترند؛ و هرچه $f_\theta$ باریک‌تر باشد، به اطلاعات پیشین خود مطمئن‌تریم.
\end{note}

\begin{center}
\renewcommand{\arraystretch}{1.45}
\begin{tabular}{>{\raggedright\arraybackslash}p{3.2cm}|>{\raggedright\arraybackslash}p{5.2cm}|>{\raggedright\arraybackslash}p{5.2cm}}
& دیدگاه کلاسیک (فراوانی) & دیدگاه بیزی\\\hline
پارامتر $\theta$ & ثابت ولی نامعلوم & متغیر تصادفی با توزیع پیشین $f_\theta$\\
منبع اطلاعات & فقط داده‌ها (درست‌نمایی) & درست‌نمایی + پیشین\\
خروجی اصلی & یک عدد (مثلاً \lr{MLE}) یا بازه اطمینان & یک \textbf{توزیع} کامل: توزیع پسین\\
نمونه تخمین‌گر & $\hat\theta_\ML$ & $\hat\theta_{\mathrm{MAP}}$، میانگین پسین
\end{tabular}
\end{center}

\subsection{سناریو تخمین بیزی: از توزیع پیشین به سمت توزیع پسین}
\begin{itemize}
\item فرض کنید متغیر تصادفی $X$ دارای توزیع پارامتریزه‌شده $f_X(x;\theta)$ است، که پارامتر(های) $\theta$ اعضای یک فضای مشخص پارامترها به نام $\Theta$ هستند.
\item پارامتر(های) $\theta$ خود به صورت یک متغیر تصادفی با توزیع مفروض و دانسته‌شده $f_\theta(\theta)$ مدل‌سازی شده است. این توزیع در بخش‌هایی که به صورت \textbf{پیشین \lr{(a priori)}} می‌دانیم محتمل‌تر هستند بزرگ‌تر و در سایر بخش‌ها کوچک‌تر است.
\begin{itemize}
\item میزان این بزرگی و کوچکی در بخش‌های مختلف از فضای پارامترها (یعنی $\Theta$) به میزان \textbf{اطمینان} ما از این اطلاعات پیشین بستگی دارد.
\end{itemize}
\item همچنین، نمونه‌های $X_1,\ldots,X_n$ به صورت \lr{i.i.d.} از توزیع به دست آمده‌اند. این مشاهدات دارای توزیع شرطی زیر هستند که همان \textbf{درست‌نمایی \lr{(Likelihood)}} در تخمین‌گر \lr{MLE} محسوب می‌شود:
\[
f_{X_{1:n}|\theta}\big(X_1,\ldots,X_n\given\theta\big)=\prod_{i=1}^nf_X\big(X_i\given\theta\big).
\]
\end{itemize}
نماد $f_X(x;\theta)$ در دیدگاه کلاسیک (که $\theta$ فقط یک «برچسب» برای انتخاب توزیع بود) اینجا به $f_X(x\given\theta)$ تبدیل شده است: چگالی \textbf{شرطی} $X$ به شرط اینکه متغیر تصادفی $\theta$ مقدار مشخصی گرفته باشد.

\begin{warning}
\warntitle در مدل بیزی، نمونه‌ها \textbf{به شرط $\theta$} مستقل‌اند (\lr{conditionally i.i.d.})، اما بدون شرط مستقل \textbf{نیستند}! چون همه آن‌ها به یک $\theta$ مشترک وابسته‌اند، دیدن $X_1$ درباره $\theta$ اطلاعات می‌دهد و در نتیجه پیش‌بینی ما از $X_2$ را عوض می‌کند. مثلاً اگر سکه‌ای با احتمال شیر نامعلوم ۱۰ بار پشت سر هم شیر بیاید، احتمال شیر آمدن بار یازدهم را (بدون شرط) بیشتر از قبل تخمین می‌زنیم.
\end{warning}

\subsection{قانون بیز برای پارامترها}
با استفاده از این اطلاعات و نحوه مدل‌سازی، و با بهره‌گیری از قانون بیز، می‌توان به یک توزیع \textbf{پسین \lr{(a posteriori)}} برای پارامتر(های) $\theta$ رسید. گام به گام: طبق تعریف چگالی شرطی، چگالی توأم $(\theta,X_{1:n})$ را به دو شکل می‌نویسیم:
\[
f_{\theta,X_{1:n}}\big(\theta,X_{1:n}\big)=f_{X_{1:n}|\theta}\big(X_{1:n}\given\theta\big)\,f_\theta(\theta)=f_{\theta|X_{1:n}}\big(\theta\given X_{1:n}\big)\,f_{X_{1:n}}\big(X_{1:n}\big).
\]
با تقسیم بر $f_{X_{1:n}}(X_{1:n})$ و جایگذاری حاصل‌ضرب درست‌نمایی:
\[
\boxed{\ f_{\theta|X_{1:n}}\big(\theta\given X_1,\ldots,X_n\big)=\frac{\prod_{i=1}^nf_X\big(X_i\given\theta\big)\cdot f_\theta(\theta)}{f_{\underline X}\big(X_1,\ldots,X_n\big)}\ }
\]
مخرج، چگالی حاشیه‌ای داده‌هاست که با قانون احتمال کل (انتگرال‌گیری روی همه مقادیر پارامتر) به دست می‌آید:
\[
f_{\underline X}\big(X_1,\ldots,X_n\big)=\int_\Theta\prod_{i=1}^nf_X\big(X_i\given\theta'\big)\,f_\theta(\theta')\dd\theta' .
\]
این مخرج \textbf{به $\theta$ وابسته نیست} (روی $\theta'$ انتگرال گرفته شده)، و تنها نقشش این است که مساحت زیر توزیع پسین را برابر ۱ کند. به همین دلیل معمولاً می‌نویسیم:
\[
f_{\theta|X_{1:n}}\big(\theta\given X_1,\ldots,X_n\big)\ \propto\ \prod_{i=1}^nf_X\big(X_i\given\theta\big)\cdot f_\theta(\theta),
\]
یا به زبان کلامی:
\[
\text{پسین}\ \propto\ \text{درست‌نمایی}\times\text{پیشین}.
\]

\begin{center}
\begin{tikzpicture}
\end{tikzpicture}\\
{\small پیشین، درست‌نمایی و پسین برای احتمال شیر آمدن یک سکه: پیشین حول $0.5$ متمرکز است (سکه‌ها معمولاً سالم‌اند) و داده‌ها ($7$ شیر در $10$ پرتاب) حول $0.7$. پسین حاصل‌ضرب این دو (پس از نرمال‌سازی) است و \textbf{بین} آن‌ها قرار می‌گیرد؛ همچنین از هر دو باریک‌تر است، چون اطلاعات دو منبع روی هم جمع شده است.}
\end{center}

\subsection{مدل تولید داده بیزی}
\textbf{نکته تکمیلی.} مدل خاص فرض‌شده برای تولید مشاهدات در سناریوهای تخمین بیزی \lr{(Bayesian Data Generation Model)} به شکل زیر است:
\[
\boxed{\ \begin{aligned}
&\theta\ \sim\ f_\theta(\theta)\\
&X_1,\ldots,X_n\ \sim\ f_{X|\theta}\big(X\given\theta\big)\quad(\text{\lr{i.i.d.}}\ \text{به شرط}\ \theta)
\end{aligned}\ }
\]
یعنی داده‌ها در \textbf{دو مرحله} تولید می‌شوند: ابتدا یک مقدار پارامتر از توزیع پیشین قرعه‌کشی می‌شود، سپس با این پارامتر ثابت، $n$ نمونه مستقل تولید می‌شود. ما فقط مرحله دوم را می‌بینیم و می‌خواهیم مرحله اول را «حدس بزنیم». (مثال: ابتدا از یک کیسه سکه‌های مختلف، یکی به تصادف برداشته می‌شود؛ سپس همان سکه $n$ بار پرتاب می‌شود.)

\subsection{از توزیع پسین به تخمین نقطه‌ای}
توزیع پسین همه چیزهایی را که (با ترکیب داده و پیشین) درباره $\theta$ می‌دانیم خلاصه می‌کند. در مرحله آخر، باید هدف اصلی برآورده شود: شکلی از تخمین برای پارامتر $\theta$ برحسب توزیع پسین به دست آمده:
\begin{itemize}
\item \textbf{محاسبه میانگین آن} به عنوان یک تخمین نقطه‌ای برای $\theta$ \lr{(Mean Posterior Estimation)}:
\[
\hat\theta_{\mathrm{PM}}=\E\big[\theta\given X_{1:n}\big]=\int_\Theta\theta\,f_{\theta|X_{1:n}}\big(\theta\given X_{1:n}\big)\dd\theta .
\]
\item \textbf{محاسبه بیشینه تابع توزیع پسین} \lr{(Maximum A Posteriori یا MAP)}، یعنی محتمل‌ترین مقدار (مُد) پسین.
\item و در صورتی که اعمال بالا سخت و غیرممکن می‌نمود، \textbf{نمونه‌گیری \lr{(sampling)}} از توزیع پسین و تخمین میانگین یا مد آن از روی نمونه‌ها. به این‌گونه روش‌ها در اصطلاح \lr{Markov-Chain Monte-Carlo} یا \lr{MCMC} می‌گویند. (ایده: حتی وقتی مخرج قانون بیز قابل محاسبه نیست، می‌توان زنجیره‌ای تصادفی ساخت که فقط با داشتن «درست‌نمایی × پیشین» نمونه‌هایی از پسین تولید کند؛ سپس میانگین نمونه‌ها، میانگین پسین را تقریب می‌زند.)
\end{itemize}

\begin{note}
\notetitle \textbf{چرا میانگین پسین؟} میانگین پسین، خطای مربعات مورد انتظار (به شرط داده‌ها) را کمینه می‌کند. برای هر عدد $c$ (که می‌تواند تابعی از داده‌ها باشد)، با اضافه و کم کردن $m=\E[\theta\given X_{1:n}]$:
\[
\E\big[(\theta-c)^2\given X_{1:n}\big]=\E\big[(\theta-m)^2\given X_{1:n}\big]+2(m-c)\underbrace{\E\big[\theta-m\given X_{1:n}\big]}_{=0}+(m-c)^2=\Var\big(\theta\given X_{1:n}\big)+(m-c)^2,
\]
که کمینه آن در $c=m$ است. به همین دلیل این تخمین‌گر را \lr{MMSE} (کمینه میانگین مربعات خطا) هم می‌نامند. در مقابل، \lr{MAP} «محتمل‌ترین» مقدار را برمی‌گرداند و معمولاً محاسبه آسان‌تری دارد (نیازی به مخرج قانون بیز و انتگرال ندارد).
\end{note}

\section{تخمین‌گر بیشینه توزیع پسین \lr{(MAP)}}

\begin{definition}[تخمین‌گر \lr{Maximum A Posteriori (MAP)}]
در این روش تخمین، مقداری از پارامتر $\theta$ (نقطه‌ای از فضای پارامتر $\Theta$) به عنوان تخمین‌گر بازگردانده می‌شود که توزیع پسین \lr{(posterior distribution)} را بیشینه نماید:
\[
\hat\theta_{\mathrm{MAP}}\triangleq\argmax_{\theta\in\Theta}\ \frac{\prod_{i=1}^nf\big(X_i\given\theta\big)f(\theta)}{f(X)}\ =\ \argmax_{\theta\in\Theta}\ \prod_{i=1}^nf\big(X_i\given\theta\big)f(\theta),
\]
یا معادلاً:
\[
\hat\theta_{\mathrm{MAP}}\triangleq\argmax_{\theta\in\Theta}\ \sum_{i=1}^n\log f\big(X_i\given\theta\big)+\log f(\theta).
\]
\end{definition}

\textbf{توضیح دو تساوی:}
\begin{itemize}
\item \textbf{حذف مخرج:} $f(X)=f_{\underline X}(X_1,\ldots,X_n)$ به $\theta$ وابسته نیست و مثبت است. ضرب یا تقسیم تابع هدف در یک عدد مثبت ثابت، محل بیشینه را تغییر نمی‌دهد (فقط مقدار بیشینه را تغییر می‌دهد). پس برای یافتن \lr{MAP} اصلاً نیازی به محاسبه انتگرال مخرج نیست؛ این یک مزیت محاسباتی بزرگ است.
\item \textbf{لگاریتم:} $\log$ تابعی اکیداً صعودی است، پس محل بیشینه را حفظ می‌کند؛ و لگاریتم حاصل‌ضرب، مجموع لگاریتم‌هاست. همان ترفند \lr{MLE} (جلسه ۲۱).
\end{itemize}

\textbf{مقایسه با \lr{MLE}:}
\[
\hat\theta_\ML=\argmax_\theta\ \underbrace{\sum_{i=1}^n\log f(X_i\given\theta)}_{\text{\lr{log-likelihood}}},
\qquad
\hat\theta_{\mathrm{MAP}}=\argmax_\theta\ \Big[\underbrace{\sum_{i=1}^n\log f(X_i\given\theta)}_{\text{\lr{log-likelihood}}}+\underbrace{\log f(\theta)}_{\text{\lr{log-prior}}}\Big].
\]
تنها تفاوت، جمله $\log f(\theta)$ است. نکته شهودی مهم: جمله درست‌نمایی مجموع $n$ جمله است و با $n$ رشد می‌کند، اما جمله پیشین یک جمله ثابت است. پس:
\begin{itemize}
\item برای $n$ کوچک، پیشین نقش مهمی دارد و تخمین را به سمت مقادیر محتمل پیشین می‌کشد؛
\item برای $n$ بزرگ، درست‌نمایی غالب می‌شود و $\hat\theta_{\mathrm{MAP}}\approx\hat\theta_\ML$.
\end{itemize}
این دقیقاً همان رفتاری است که در مثال جراحان به طور شهودی انتظار داشتیم! (نکات بیشتر درباره \lr{MAP}، از جمله نقش پیشین به عنوان «جمله تنظیم‌کننده»، و مثال گاوسی در جلسه ۲۷ می‌آید.)

\begin{example}[سکه سالم یا سکه دست‌کاری‌شده؟]
در یک کیسه، ۹۰٪ سکه‌ها سالم ($\theta=0.5$) و ۱۰٪ دست‌کاری‌شده با احتمال شیر $\theta=0.8$ هستند. سکه‌ای به تصادف از کیسه برداشته و $5$ بار پرتاب می‌شود و نتیجه $4$ شیر و $1$ خط است. (الف) \lr{MLE} و \lr{MAP} را برای $\theta$ بیابید. (ب) میانگین پسین را حساب کنید. (ج) اگر به جای آن $20$ پرتاب با $16$ شیر دیده بودیم، \lr{MAP} چه می‌شد؟
\begin{solution}
فضای پارامتر گسسته است: $\Theta=\{0.5,\ 0.8\}$ و پیشین $\Prob(\theta=0.5)=0.9$، $\Prob(\theta=0.8)=0.1$ (اینجا به جای چگالی، تابع جرم احتمال داریم).

\textbf{درست‌نمایی} یک دنباله مشخص با $4$ شیر و $1$ خط برابر $\theta^4(1-\theta)$ است:
\[
L(0.5)=0.5^5=0.03125,\qquad L(0.8)=0.8^4\times0.2=0.08192 .
\]
(الف) \textbf{\lr{MLE}:} چون $L(0.8)>L(0.5)$، داریم $\hat\theta_\ML=0.8$.

\textbf{\lr{MAP}:} حاصل‌ضرب درست‌نمایی در پیشین:
\[
L(0.5)\times0.9=0.028125,\qquad L(0.8)\times0.1=0.008192 .
\]
بزرگ‌تر مربوط به $\theta=0.5$ است، پس $\hat\theta_{\mathrm{MAP}}=0.5$. شواهد داده‌ها (نسبت درست‌نمایی $0.08192/0.03125\approx2.6$ به نفع سکه دست‌کاری‌شده) برای غلبه بر پیشین (نسبت $9$ به نفع سکه سالم) کافی نیست.

(ب) برای پسین کامل، نرمال می‌کنیم: مجموع $0.028125+0.008192=0.036317$ (همان مخرج قانون بیز):
\[
\Prob(\theta=0.5\given D)=\frac{0.028125}{0.036317}=0.774,\qquad \Prob(\theta=0.8\given D)=\frac{0.008192}{0.036317}=0.226 .
\]
میانگین پسین: $0.5\times0.774+0.8\times0.226=0.568$. (این عدد حتی عضو $\Theta$ نیست! میانگین پسین لزوماً یک مقدار ممکن پارامتر نیست.)

(ج) با $16$ شیر در $20$ پرتاب، نسبت درست‌نمایی:
\[
\frac{L(0.8)}{L(0.5)}=\frac{0.8^{16}\times0.2^4}{0.5^{20}}\approx47.2,
\]
و نسبت پسین $47.2\times\frac{0.1}{0.9}\approx5.25$ به نفع $\theta=0.8$؛ یعنی $\Prob(\theta=0.8\given D)\approx0.84$ و $\hat\theta_{\mathrm{MAP}}=0.8$. با داده بیشتر، شواهد بر پیشین غلبه کرد.
\end{solution}
\end{example}

\begin{note}
\notetitle \textbf{بیز به زبان نسبت‌ها.} در مثال بالا از یک فرم بسیار مفید قانون بیز استفاده کردیم: برای دو مقدار $\theta_1,\theta_2$،
\[
\frac{\Prob(\theta_1\given D)}{\Prob(\theta_2\given D)}=\frac{\Prob(D\given\theta_1)}{\Prob(D\given\theta_2)}\times\frac{\Prob(\theta_1)}{\Prob(\theta_2)},
\]
یعنی «نسبت پسین = نسبت درست‌نمایی × نسبت پیشین». مخرج قانون بیز در این نسبت حذف می‌شود.
\end{note}

\section{تمرین‌های تکمیلی}

\begin{exercise}[تعداد نمونه لازم: چبی‌شف در برابر حد مرکزی]
می‌خواهیم میانگین یک توزیع با واریانس $\sigma^2$ را با دقت $\pm0.1\sigma$ و اطمینان ۹۵٪ تخمین بزنیم. تعداد نمونه لازم را (الف) با بازه چبی‌شف و (ب) با بازه تقریبی حد مرکزی حساب کنید.
\end{exercise}
\begin{solution}
(الف) باید $\frac{\sigma}{\sqrt{n\delta}}\le0.1\sigma$، یعنی $n\ge\frac{1}{0.01\times\delta}=\frac{1}{0.01\times0.05}=2000$.

(ب) باید $1.96\frac{\sigma}{\sqrt n}\le0.1\sigma$، یعنی $n\ge19.6^2=384.2$، پس $n=385$.

چبی‌شف حدود ۵ برابر نمونه بیشتر می‌خواهد ($\frac{4.472^2}{1.96^2}\approx5.2$)؛ این بهای تضمین بدون فرض است.
\end{solution}

\begin{exercise}[بازه چبی‌شف برای نسبت]
در نظرسنجی جلسه قبل ($n=400$، $\hat p=0.45$) بازه اطمینانی با سطح حداقل ۹۵٪ با استفاده از نامساوی چبی‌شف بسازید.
\end{exercise}
\begin{solution}
برای متغیر برنولی $\Var X=p(1-p)\le\frac14$، پس $\sigma=\frac12$ کران معتبری است:
\[
\frac{\sigma}{\sqrt{n\delta}}=\frac{0.5}{\sqrt{400\times0.05}}=\frac{0.5}{\sqrt{20}}=0.112
\quad\Rightarrow\quad[0.338,\ 0.562].
\]
در مقایسه با بازه تقریبی $[0.401,\ 0.499]$ (نصف طول $0.049$)، بیش از دو برابر پهن‌تر است.
\end{solution}

\begin{exercise}[سه سکه]
در جعبه‌ای سه سکه با احتمال شیر $0.2$، $0.5$ و $0.8$ هست. یکی به تصادف (با احتمال برابر) انتخاب و سه بار پرتاب می‌شود: شیر، شیر، خط. توزیع پسین، \lr{MAP}، \lr{MLE} و میانگین پسین را بیابید.
\end{exercise}
\begin{solution}
درست‌نمایی $\theta^2(1-\theta)$: $L(0.2)=0.032$، $L(0.5)=0.125$، $L(0.8)=0.128$. پیشین یکنواخت ($\frac13$ برای هر سکه) در نسبت‌ها حذف می‌شود؛ پس پسین متناسب با درست‌نمایی است. مجموع $0.285$:
\[
\Prob(0.2\given D)=0.112,\qquad\Prob(0.5\given D)=0.439,\qquad\Prob(0.8\given D)=0.449 .
\]
$\hat\theta_{\mathrm{MAP}}=\hat\theta_\ML=0.8$ (با پیشین یکنواخت، \lr{MAP} و \lr{MLE} یکی‌اند؛ جلسه ۲۷). میانگین پسین: $0.2(0.112)+0.5(0.439)+0.8(0.449)=0.601$. توجه کنید که \lr{MAP} به سختی $0.8$ را بر $0.5$ ترجیح داده است؛ میانگین پسین این عدم قطعیت را بهتر نشان می‌دهد.
\end{solution}

\begin{exercise}[\lr{MAP} برای پواسون با پیشین نمایی]
نمونه‌های $X_1,\ldots,X_n$ به شرط $\lambda$ به صورت \lr{i.i.d.} از $\mathrm{Poisson}(\lambda)$ هستند و پیشین $\lambda\sim\mathrm{Exp}(\beta)$، یعنی $f(\lambda)=\beta e^{-\beta\lambda}$ برای $\lambda\ge0$. تخمین \lr{MAP} را بیابید و با \lr{MLE} مقایسه کنید.
\end{exercise}
\begin{solution}
$\log f(X_i\given\lambda)=X_i\log\lambda-\lambda-\log X_i!$ و $\log f(\lambda)=\log\beta-\beta\lambda$. تابع هدف (با حذف ثابت‌ها):
\[
g(\lambda)=\Big(\sum_iX_i\Big)\log\lambda-n\lambda-\beta\lambda .
\]
مشتق: $g'(\lambda)=\frac{\sum_iX_i}{\lambda}-(n+\beta)=0$، پس
\[
\hat\lambda_{\mathrm{MAP}}=\frac{\sum_{i=1}^nX_i}{n+\beta},
\]
و چون $g''(\lambda)=-\frac{\sum_iX_i}{\lambda^2}<0$، بیشینه است. در مقایسه با $\hat\lambda_\ML=\frac{\sum_iX_i}{n}$، مخرج به اندازه $\beta$ بزرگ‌تر است: پیشین (که مقادیر کوچک $\lambda$ را محتمل‌تر می‌داند) تخمین را به سمت صفر «جمع» کرده است. هرچه $\beta$ بزرگ‌تر (پیشین مطمئن‌تر)، این اثر بیشتر؛ و وقتی $n\gg\beta$، دو تخمین تقریباً برابرند.
\end{solution}

\begin{exercise}[به‌روزرسانی ترتیبی]
داده‌ها در دو مرحله $D_1$ و $D_2$ می‌رسند که به شرط $\theta$ مستقل‌اند. نشان دهید اگر پسین مرحله اول را به عنوان پیشین مرحله دوم به کار ببریم، همان پسینی به دست می‌آید که با پردازش یکجای همه داده‌ها حاصل می‌شود.
\end{exercise}
\begin{solution}
پردازش یکجا: $f(\theta\given D_1,D_2)\propto f(D_1,D_2\given\theta)f(\theta)=f(D_2\given\theta)\,f(D_1\given\theta)f(\theta)$ (استقلال شرطی). پردازش ترتیبی: پس از مرحله اول $f(\theta\given D_1)\propto f(D_1\given\theta)f(\theta)$، و در مرحله دوم $f(D_2\given\theta)\,f(\theta\given D_1)\propto f(D_2\given\theta)f(D_1\given\theta)f(\theta)$. دو عبارت متناسب‌اند و هر دو چگالی‌اند (مساحت ۱)، پس برابرند. شهود: «پسین امروز، پیشین فرداست»؛ یادگیری بیزی می‌تواند به صورت برخط \lr{(online)} و داده به داده انجام شود.
\end{solution}

\begin{exercise}[میانه پسین]
نشان دهید عددی که خطای قدرمطلق مورد انتظار پسین، یعنی $\E\big[|\theta-c|\given D\big]$، را کمینه می‌کند، \textbf{میانه} توزیع پسین است.
\end{exercise}
\begin{solution}
با $F$ تابع توزیع تجمعی پسین: $h(c)=\E|\theta-c|=\int_{-\infty}^c(c-t)f(t)\dd t+\int_c^\infty(t-c)f(t)\dd t$. با مشتق‌گیری (قاعده لایب‌نیتس؛ جمله‌های مرزی صفرند):
\[
h'(c)=\int_{-\infty}^cf(t)\dd t-\int_c^\infty f(t)\dd t=F(c)-\big(1-F(c)\big)=2F(c)-1 .
\]
این مشتق برای $F(c)<\frac12$ منفی و برای $F(c)>\frac12$ مثبت است؛ پس کمینه در $F(c)=\frac12$، یعنی میانه. (مشابه تمرین توزیع لاپلاس در جلسه ۲۲.) خلاصه: خطای مربعی $\leftarrow$ میانگین پسین، خطای قدرمطلق $\leftarrow$ میانه پسین، و \lr{MAP} $\leftarrow$ مُد پسین.
\end{solution}
